\section{Multiple inputs: a declared population, not arbitrary waves}
\label{sec:data}

Begin with $L=160$ and $\epsilon=.025$. A training example contains the
entire complex initial profile, the actual local propagation increment,
and its reference output. Future reference fields are training targets
only. The same trained operator must accept a new initial profile
without optimization. If adaptation is allowed, label and time it
separately.

\subsection{Three complementary input families}
The following ranges are proposed starting values for a bounded pilot,
not ranges established by a physical-validity theorem. Validate a small
batch before freezing them. Record all rejected draws and their reason;
an unreported rejection rule changes the population being learned.

\paragraph{Localized packets: focusing and asymmetry.}
Use a periodized version of
\[
u_0(\tau)=a\exp\left[-\frac{1-ic}{2\sigma^2}(\tau-\tau_0)^2+i\phi\right],
\quad
a\in[.5,1],\quad \sigma\in[6,12],\quad c\in[-.5,.5].
\]
Vary position and phase as well. Periodize the complex function by summing
translated copies, with a documented derivative-level truncation tolerance.
A cropped Gaussian is not exactly periodic. This family tests amplitude,
width, chirp, focusing distance and asymmetric deformation. Hold out
entire parameter draws, not propagation stations. Sech-shaped packets
and two separated packets can later be withheld \emph{families}, not
silently mixed into an interpolation score.

\paragraph{Modulated wavetrains: instability and recurrence.}
Use
\[
u_0(\tau)=a e^{i\phi}
 \left[1+b_+e^{ik_m\tau}+b_-e^{-ik_m\tau}\right].
\]
Choose $a\in[.5,1]$, integer modes spanning stable, near-boundary and
unstable portions of the matrix in Appendix~\ref{app:mi}, and independent
sideband phases. Use $|b_\pm|$ between $.01$ and $.05$ for a nonlinear
pilot; use much smaller paired disturbances for the linearization test.
These waves intentionally fill the periodic domain. They do not need a
localized-packet edge criterion.

\paragraph{Smooth random profiles: functional variation.}
Sample complex Fourier coefficients with a prescribed decaying spectrum,
for example a finite series with $|m|\leq16$ and weights
$\exp[-(m/m_c)^2]$, $m_c\in[4,10]$, followed by a declared RMS-amplitude
choice. Coefficients at $m$ and $-m$ are independent. Bound both peak
amplitude and derivative-weighted bandwidth by rejection or a stated
sampling rule. This is a finite, smooth input distribution with variable
phase, not a claim to represent all sea states. Record its seed and
coefficient vector. An alternative spectral law is a separate
distribution-shift test.

\paragraph{Information that cannot be discarded.}
Intensity $|u_0|^2$ alone does not specify the initial-value problem:
opposite chirps can have the same intensity and different evolution.
Sparse or noisy measurements require an observation model, a
reconstruction method and uncertainty evaluation. They are a later
inverse problem; they must not be presented as another complete-field
input channel with missing phase filled in from the reference.

\subsection{Splits, grids and provenance}
Store reference fields as complex128 arrays $[N_\xi,N_\tau]$ with
$\tau_j=-L/2+jL/N$, excluding the duplicated endpoint. Store explicit
propagation coordinates. The existing runner's rounded output indices
can create unequal increments; validate the actual spacing or condition
the operator on each actual increment.

Assign one immutable \emph{split group} to each base initial draw.
All its windows, refinements, translated or phase-rotated copies,
paired disturbances, and coefficient variants remain in the same split.
Test for duplicate field hashes as well as mislabeled groups.
Generate parameter ranges and seed lists before fitting. Fit global
normalization on training fields only. Use validation for model selection;
open the final test set once the protocol and checkpoint are frozen.
Snapshots from one field are correlated samples, not independent cases.

The current metadata validator is insufficient by itself. The new loader
must check actual shapes, dtypes, finite values, coordinates, hashes,
model ID, epsilon, grid origin, and qualification status. Preserve case
schema 2 and add a separate dataset manifest for group and split records.

Downsample through Fourier restriction, not viewer decimation. Measure
projection and reconstruction errors in $L^2,H^1,H^2,H^3$ over the
qualified reference interval. A 512- or 1024-point learning grid is only a
candidate; require its representation error to consume at most one tenth
of the intended surrogate error budget. Choose the number of learned
Fourier modes through a small validation sweep, such as $32,64,128$,
after inspecting the physical spectra. These numbers do not specify the
FNO output bandwidth.

\subsection{A campaign that can stop early}
Use equal numbers from the three families initially.
\begin{center}
\small
\begin{tabular}{p{.28\linewidth}rrp{.25\linewidth}}
\toprule
Stage & Groups & Train/val/test & Purpose\\
\midrule
Execution smoke & 12 & $6/3/3$ & Load, train, resume and evaluate\\
Bounded pilot & 48 & $36/6/6$ & Select viable method and tolerances\\
Main comparison, if justified & 384 & $288/48/48$ & Estimate held-out performance\\
\bottomrule
\end{tabular}
\end{center}
These are proposed budgets, not power calculations or guaranteed
sufficient sample sizes. Six pilot test groups cannot substantiate broad
generalization. Use learning curves on nested training subsets to decide
whether more groups are useful. Keep the test population fixed while
increasing training data.

At $N=4096$ and 81 saved stations on $[0,4]$, one complex128 trajectory
requires about $5.06$ MiB before compression. Thus 48 and 384 base
trajectories require about 243 MiB and 1.90 GiB, respectively, excluding
refinements and checkpoints. A saved interval $.05$ is a proposed
\emph{learning} increment; the reference integration step $.00125$
is a separate quantity. Extra stations are needed for continuous-distance
training and near fast events.

Qualify representative extremes first, then screen every generated
trajectory; refine any case whose derivative tails or defects approach
the gate. Use full three-level grid and step comparisons on the sentinel
panel and difficult cases, and at least an independent refined
comparison for accepted evaluation cases. Do not assume the three old
demonstrations qualify a whole new distribution.

For the main campaign's first long-distance evaluation, preselect 12 held-out groups
(four per family) by input parameters alone. Extend $4\to8\to16$ and
only then, if useful and qualified, to 32. At 641 stations a 12-case
$N=4096$ panel needs about 481 MiB of base fields. The same panel supports
paired comparisons across methods; do not repeatedly select its easiest
members. Benchmark actual generation and training costs before launching
the main campaign.

Vary epsilon only after the fixed-parameter question works. A proposed
second stage uses $.015,.025,.04$ in training and disjoint groups at
$.02,.03$ for interpolation. Keep $.05$ a separately labeled
extrapolation probe, subject to new qualification and physical screens.
Treat $\epsilon=0$ as an analytic/numerical control. Changing $L$ is a
different domain, and changing transverse dimensionality or dissipation
is a different model.
